\chapter{Вильсоновский селектор трёх Majorana-core мод}

Полярный аудит ordinary \(B-L\)-ветви исправил важный подсчёт: если vortex
поколенчески-слеп, то до дополнительных family-взаимодействий он создаёт
три вещественные core-моды
\begin{equation}
 \chi=(\chi_1,\chi_2,\chi_3)^T\in\mathbb R^3_{gen},
\end{equation}
а не одну. Это закрывает прежнюю попытку получить \(24-1\) одним лишь
проектированием pairing-матрицы, но открывает механизм, уже присутствующий
в части II.B.

\section{Антисимметричная связь на core}

Непрерывная вильсоновская ветвь построила правильное вращение
\(R_n(\theta_*)\in SO(3)\) на family triplet, где
\begin{equation}
 \cos\theta_*=\frac{26-9\sqrt{15}}{11}.
\end{equation}
Его главный генератор имеет вид
\begin{equation}
 K_n=\operatorname{Log}R_n(\theta_*)
     =\theta_*[n]_\times,
 \qquad
 ([n]_\times)_{ab}=\epsilon_{abc}n_c.
\end{equation}
Он вещественен, антисимметричен и при \(\theta_*\ne0\) имеет
\begin{equation}
 \operatorname{rank}K_n=2,
 \qquad
 \ker K_n=\mathbb R n.
\end{equation}

Поэтому Majorana boundary term
\begin{equation}
 S_{\rm core,fam}
 =\frac{i}{2}\int_\gamma \chi^T K_n\chi
\end{equation}
спаривает две из трёх core-мод в одну ненулевую fermionic пару и оставляет
ровно одну вещественную нулевую моду вдоль оси \(n\). Это не требует
rank-one Yukawa matrix \(uu^T\): selector возникает как kernel уже
существующего ориентированного family generator.

\section{Независимость ранга от выбора вакуума}

Factor-axis functional сокращает восемь совместных осей до двух
симметрийно-связанных минимумов. Машинная проверка для обеих осей даёт
\begin{equation}
 \Spec(iK_n)=\{-\theta_*,0,+\theta_*\},
 \qquad
 \dim\ker K_n=1.
\end{equation}
Замена ориентации \(n\mapsto-n\) меняет \(K_n\mapsto-K_n\), но сохраняет
kernel projector \(P_n=nn^T\). Следовательно, знак ориентированного
вакуумного выбора не влияет на число surviving Majorana modes.

\begin{proposition}[Условное снятие generation multiplicity]
Если family Wilson connection ограничивается на defect core именно через
\(K_n\), то поколенчески-слепой vortex сначала создаёт три локальные
Majorana-моды, после чего антисимметричная family-связь оставляет одну:
\begin{equation}
 3\longrightarrow1.
\end{equation}
При наличии ambient-to-core restriction map это восстанавливает кандидат на
rank-one kernel и complement rank \(24-1=23\) без ручной вставки \(uu^T\).
\end{proposition}

\section{Оставшийся parent-action gate}

Положительный алгебраический результат ещё не является полным замыканием.
Текущее действие не выводит в одном функционале:
\begin{enumerate}
 \item ограничение той же \(SO(3)\) family connection на \(B-L\) defect core;
 \item коэффициент и знак \(S_{\rm core,fam}\) вместе с vortex action;
 \item ambient-to-core map из \(\mathbb R^{24}\) в три локальные моды;
 \item collective stiffness, превращающую complement rank в
 \(23+\pi^{-1}\), а не только в кратность determinant.
\end{enumerate}

\begin{nogocriterion}[Wilson--defect bridge gate]
Если \(S_{\rm core,fam}\) добавляется как независимый boundary term, не
следующий из того же graded action, что и root vortex и Wilson gap-functional,
то rank-one selector имеет статус конструктивного механизма, но не вывода
S2T. Положительное закрытие требует одного covariant derivative или одного
boundary superconnection, чья редукция одновременно даёт vortex BdG operator
и \(K_n\)-связь на его трёхмерном kernel bundle.
\end{nogocriterion}

Таким образом, проблема больше не состоит в отсутствии возможного generation
selector. Конкретный selector найден внутри уже построенной Wilson-ветви.
Открытым остаётся мост между family Wilson connection и \(B-L\) defect
parent action.

Машинный аудит сохранён в
\texttt{s2t\_family\_wilson\_majorana\_core\_selector\_results.json}.